% 编译方式：XeLaTeX
\documentclass[UTF8,11pt,a4paper,fontset=none]{ctexart}

\usepackage[a4paper,top=1.7cm,bottom=1.8cm,left=1.75cm,right=1.75cm]{geometry}
\usepackage[table]{xcolor}
\usepackage{booktabs}
\usepackage{tabularx}
\usepackage{longtable}
\usepackage{array}
\usepackage{enumitem}
\usepackage{tcolorbox}
\tcbuselibrary{skins,breakable}
\usepackage{hyperref}
\usepackage{fancyhdr}
\usepackage{titlesec}
\usepackage{setspace}
\usepackage{tikz}
\usepackage{fontspec}
\usetikzlibrary{arrows.meta,positioning,fit}

\setCJKmainfont{Noto Serif CJK SC}
\setCJKsansfont{Noto Sans CJK SC}
\setCJKmonofont{Noto Sans Mono CJK SC}
\setmainfont{Noto Serif CJK SC}
\setsansfont{Noto Sans CJK SC}
\setmonofont{Noto Sans Mono CJK SC}

\definecolor{navy}{HTML}{102A43}
\definecolor{ocean}{HTML}{1769AA}
\definecolor{teal}{HTML}{0F8B8D}
\definecolor{mint}{HTML}{E6F4F1}
\definecolor{sky}{HTML}{EAF3FA}
\definecolor{slate}{HTML}{486581}
\definecolor{ink}{HTML}{243B53}
\definecolor{muted}{HTML}{627D98}
\definecolor{line}{HTML}{D9E2EC}
\definecolor{paper}{HTML}{F7FAFC}
\definecolor{amber}{HTML}{D97706}

\hypersetup{
  colorlinks=true,
  linkcolor=ocean,
  urlcolor=teal,
  pdfauthor={赛先生品牌内容系统},
  pdftitle={第一部分建设成果报告：权威源自动采集与 AI 证据入库}
}

\setlength{\parindent}{2em}
\setlength{\parskip}{0.3em}
\setlength{\headheight}{22.2pt}
\setlength{\tabcolsep}{5pt}
\onehalfspacing
\raggedbottom

\pagestyle{fancy}
\fancyhf{}
\lhead{\small\sffamily\color{slate} 第一部分建设成果报告 · 权威证据采集}
\rhead{\small\sffamily\color{slate} 赛先生品牌内容系统 · 2026-07-29}
\cfoot{\small\sffamily\color{muted} \thepage}
\renewcommand{\headrulewidth}{0.5pt}
\renewcommand{\headrule}{\hbox to\headwidth{\color{ocean!65}\leaders\hrule height \headrulewidth\hfill}}

\titleformat{\section}
  {\Large\bfseries\sffamily\color{navy}}
  {\tikz[baseline=(n.base)]\node[fill=ocean,text=white,rounded corners=2pt,
    inner xsep=7pt,inner ysep=2pt,font=\bfseries\sffamily] (n) {\thesection};}
  {0.75em}{}
\titleformat{\subsection}{\large\bfseries\sffamily\color{ocean}}{\thesubsection}{0.65em}{}
\titlespacing*{\section}{0pt}{1.2em}{0.5em}
\titlespacing*{\subsection}{0pt}{0.85em}{0.25em}

\newcolumntype{Y}{>{\raggedright\arraybackslash}X}
\newcolumntype{L}[1]{>{\raggedright\arraybackslash}p{#1}}
\newcommand{\keyvalue}[2]{\noindent\textbf{\sffamily\color{navy}#1：}\textcolor{ink}{#2}\par}
\newcommand{\statuspill}[2]{%
  \tikz[baseline=(pill.base)]\node[fill=#1!13,text=#1,draw=#1!35,rounded corners=5pt,
  inner xsep=6pt,inner ysep=2pt,font=\scriptsize\bfseries\sffamily] (pill) {#2};}

\tcbset{enhanced,boxrule=0.55pt,arc=1.7mm,left=3.5mm,right=3.5mm,top=2.5mm,bottom=2.5mm}
\newtcolorbox{focusbox}[1][]{colback=mint,colframe=teal!55,
  borderline west={2.6pt}{0pt}{teal},breakable,#1}
\newtcolorbox{infobox}[1][]{colback=sky,colframe=ocean!38,
  borderline west={2.4pt}{0pt}{ocean},breakable,#1}
\newtcolorbox{boundarybox}[1][]{colback=paper,colframe=line,
  borderline west={2.4pt}{0pt}{amber},breakable,#1}

\begin{document}

\begin{center}
  \statuspill{teal}{FIRST CAPABILITY COMPLETE}\hspace{0.5em}
  \statuspill{ocean}{LIVE ACCEPTANCE PASSED}\par\vspace{0.8em}
  {\color{ocean}\large\bfseries\sffamily 第一部分建设成果报告}\par\vspace{0.5em}
  {\fontsize{23}{29}\selectfont\bfseries\sffamily\color{navy}
    权威源自动采集与 AI 证据入库}\par\vspace{0.45em}
  {\large\sffamily\color{slate}
    8 个固定权威来源 · AI 标题门禁 · 原文快照 · 全链路溯源}\par\vspace{0.9em}
  \tikz\draw[teal,line width=1.2pt] (0,0)--(3.2,0);
\end{center}

\keyvalue{报告日期}{2026 年 7 月 29 日}
\keyvalue{能力状态}{第一部分实现、独立审查、确定性测试与实时 8 来源验收均已完成}
\keyvalue{实时运行 ID}{\texttt{a174ed10-0ef1-4983-b81f-7f8d2bed84d2}}
\keyvalue{部署形态}{Python 3.11 + FastAPI + PostgreSQL/pgvector + MinIO + 独立 Scheduler/Worker}
\keyvalue{相关性规则}{\texttt{ai-title-v1}，不依赖大模型即可运行}

\begin{focusbox}
\textbf{\sffamily\color{navy}管理摘要}\quad
第一部分已从“网页抓取脚本”升级为可部署、可恢复、可审计的权威证据入口。
系统每天可从 8 个固定权威来源扫描近期列表，只对 AI、人工智能、机器人、具身智能、
大模型、算力、AI 芯片、脑机接口及相关政策标题请求详情页；无匹配时返回 0 条，
不使用无关政策或普通科研文章凑数。每条候选同时保存清洗正文、原文链接和不可变网页快照，
可直接供后续 LangGraph 总结/选题节点使用。
\end{focusbox}

\section{本阶段建设目标与业务价值}

本阶段解决后续 AI 内容生产的事实入口问题。它不负责摘要、选题、文案或生图，而是先回答
三个基础问题：\textbf{数据来自哪里、是否与 AI 主题相关、未来能否回到原文复核}。

\noindent\begin{tabularx}{\textwidth}{L{3.0cm}YY}
\toprule
\rowcolor{navy}
\color{white}\bfseries\sffamily 目标 &
\color{white}\bfseries\sffamily 已实现能力 &
\color{white}\bfseries\sffamily 直接价值 \\
\midrule
固定权威来源 & 8 个政府、高校、科研机构、企业与权威媒体来源版本化登记 & 避免模型自由浏览和未知来源污染 \\
\rowcolor{paper}
只收相关内容 & 标题相关性规则在详情请求前执行，支持 AI 相关政策 & 降低无关正文、调用成本和后续人工筛选成本 \\
保存完整证据 & 清洗正文、原始快照、URL、时间、哈希、解析版本与运行记录 & 文章修改或下线后仍可审计复现 \\
\rowcolor{paper}
每天稳定运行 & 独立调度器、持久 Job、租约、心跳、幂等与失败分类 & 服务重启后可恢复，单站失败不阻塞其他来源 \\
下游可直接消费 & 列表提供来源/标题/时间/链接，详情提供正文与快照 & 后续节点无需重复访问官网 \\
\bottomrule
\end{tabularx}

\section{来源范围与主题边界}

\subsection{固定的 8 个来源}

\noindent\begin{tabularx}{\textwidth}{L{1.1cm}L{4.0cm}L{3.6cm}Y}
\toprule
\rowcolor{ocean}
\color{white}\bfseries 级别 & \color{white}\bfseries 来源 &
\color{white}\bfseries 类型 & \color{white}\bfseries 主要用途 \\
\midrule
A & 中国政府网最新政策 & 政府/政策 & AI 规划、法规、标准、治理与产业支持政策 \\
\rowcolor{paper}
A & 北京师范大学新闻网 & 高校 & AI 教育、人才培养与学校合作的一手信息 \\
A & 中国科学院科研进展 & 科研机构 & 与 AI/智能系统明确相关的科研进展 \\
\rowcolor{paper}
A & 商汤科技新闻中心 & AI 企业 & 企业一手产品、模型、机器人与产业实践 \\
B & 新华网科技 & 权威媒体 & 科技趋势、产业背景与事实佐证 \\
\rowcolor{paper}
B & 光明网教育 & 权威教育媒体 & AI 教育、机器人赛事和学校实践 \\
B & 科技日报 & 科技媒体 & AI、具身智能、算力与科技产业动态 \\
\rowcolor{paper}
B & 中国新闻网教育 & 权威教育媒体 & AI 教育现象、治理问题与社会观察 \\
\bottomrule
\end{tabularx}

\subsection{AI 标题门禁}

规则版本 \texttt{ai-title-v1} 使用 Unicode NFKC、大小写折叠、空白与短横线归一化，
覆盖 AI/人工智能、AIGC、大模型、GPT/ChatGPT/DeepSeek、机器学习、算法、算力、
AI 芯片、计算机视觉、语音/NLP、机器人、具身智能、自动驾驶、无人机和脑机接口等主题。

政策词“规划、办法、标准、治理、通知、支持”等不能单独使文章通过；只有标题同时出现明确的
AI/智能技术词时，才作为 AI 政策接受。系统同时拒绝 \textit{Travel agents}、
\textit{Chemical agent}、商业信心指数 BCI、招生语境 UAS、智能体检、智能体育等歧义写法。

\begin{boundarybox}
\textbf{\sffamily\color{navy}保守原则}\quad
宁可暂时漏掉标题没有明确 AI 信号的文章，也不抓取无关内容。新术语通过发布新的相关性规则和
来源版本扩展，不能原地修改历史规则。
\end{boundarybox}

\section{系统实现与数据流}

\begin{center}
\begin{tikzpicture}[
  node distance=0.48cm,
  stage/.style={rectangle,draw=ocean!65,fill=sky,rounded corners=3pt,
    text width=2.35cm,minimum height=0.95cm,align=center,font=\small\bfseries\sffamily,text=navy},
  gate/.style={stage,draw=teal,fill=mint},
  arrow/.style={-Stealth,thick,color=teal!80!black}
]
  \node (source) [stage] {8 个来源\\近期列表};
  \node (safe) [stage,right=of source] {安全抓取\\SSRF · 限速 · 超时};
  \node (filter) [gate,right=of safe] {AI 标题门禁\\过滤后再抓详情};
  \node (snapshot) [stage,right=of filter] {MinIO 快照\\原始网页 · 哈希};
  \node (db) [stage,below=0.62cm of snapshot] {PostgreSQL\\Run · Job · 候选};
  \node (api) [gate,left=of db] {证据 API\\标题 · 链接 · 正文};
  \node (later) [stage,left=of api] {后续 LangGraph\\总结 · 选题 · 生成};
  \draw[arrow] (source)--(safe);
  \draw[arrow] (safe)--(filter);
  \draw[arrow] (filter)--(snapshot);
  \draw[arrow] (snapshot)--(db);
  \draw[arrow] (db)--(api);
  \draw[arrow] (api)--(later);
\end{tikzpicture}
\end{center}

\begin{enumerate}[leftmargin=2em,itemsep=0.25em]
  \item 列表响应先经过允许域名、公共 IP、重定向、类型、字节与超时安全检查；
  \item 连接器扫描有限的近期列表，合并图片空链接与文本链接，形成有序候选；
  \item 标题门禁过滤无关项，只有相关项才发起详情请求；
  \item 原始列表和相关详情以内容哈希写入 MinIO，不同字节不能覆盖同一对象；
  \item PostgreSQL 保存来源版本、运行、Job、候选、Observation、计数与快照引用；
  \item 后续节点通过 candidate ID 获取已保存正文，不把原文链接当作新的自由爬虫入口。
\end{enumerate}

\section{2026-07-29 实时验收结果}

\subsection{总体运行指标}

\noindent\begin{tabularx}{\textwidth}{L{3.2cm}Y L{3.2cm}Y}
\toprule
\rowcolor{navy}
\color{white}\bfseries 指标 & \color{white}\bfseries 结果 &
\color{white}\bfseries 指标 & \color{white}\bfseries 结果 \\
\midrule
Run ID & \texttt{\footnotesize a174ed10-0ef1-4983}\newline
         \texttt{\footnotesize b81f-7f8d2bed84d2} & 规则版本 & \texttt{ai-title-v1} \\
\rowcolor{paper}
来源任务 & 8/8 成功 & 失败任务 & 0 \\
新增相关候选 & 6 & 无相关内容来源 & 2 \\
\rowcolor{paper}
过滤无关标题 & 263 & 重复/未变化 & 0/0 \\
任务执行窗口 & 约 7.5 秒 & 本轮响应字节 & 993,007 Bytes \\
\rowcolor{paper}
未完成 Job & 0 & 抓取租约 & 0 \\
\bottomrule
\end{tabularx}

\begin{infobox}
本轮中国政府网扫描 100 条政策、中国科学院扫描 15 条科研进展，均未发现标题明确关联 AI 或
智能科技的内容，因此两个 Job 以 \texttt{no\_relevant\_items} 成功结束，没有抓取无关详情页。
这比“每站强行返回一条”更符合证据质量目标。
\end{infobox}

\subsection{各来源过滤与接受情况}

\noindent\begin{tabularx}{\textwidth}{L{4.2cm}L{3.0cm}L{2.5cm}Y}
\toprule
\rowcolor{ocean}
\color{white}\bfseries 来源 & \color{white}\bfseries 结果 &
\color{white}\bfseries 过滤数 & \color{white}\bfseries 入库数 \\
\midrule
北京师范大学新闻网 & 成功 & 22 & 1 \\
\rowcolor{paper}
中国科学院科研进展 & 无相关内容 & 15 & 0 \\
中国政府网最新政策 & 无相关内容 & 100 & 0 \\
\rowcolor{paper}
中国新闻网教育 & 成功 & 57 & 1 \\
光明网教育 & 成功 & 36 & 1 \\
\rowcolor{paper}
商汤科技新闻中心 & 成功 & 5 & 1 \\
科技日报 & 成功 & 7 & 1 \\
\rowcolor{paper}
新华网科技 & 成功 & 21 & 1 \\
\bottomrule
\end{tabularx}

\subsection{正式入库的 6 条相关内容}

\begin{longtable}{L{2.75cm}L{6.25cm}L{2.2cm}L{2.15cm}L{1.75cm}}
\toprule
\rowcolor{navy}
\color{white}\bfseries 来源 & \color{white}\bfseries 标题/原文 &
\color{white}\bfseries 发布时间 & \color{white}\bfseries 命中词 &
\color{white}\bfseries 正文字符 \\
\midrule
\endfirsthead
\toprule
\rowcolor{navy}
\color{white}\bfseries 来源 & \color{white}\bfseries 标题/原文 &
\color{white}\bfseries 发布时间 & \color{white}\bfseries 命中词 &
\color{white}\bfseries 正文字符 \\
\midrule
\endhead
北京师范大学 & 北京师范大学与北京航空航天大学签署关于人工智能+教育领域的合作框架协议
\newline\href{https://news.bnu.edu.cn/zx/ttgz/293edc4b328449368e6da0fc704432d8.htm}{查看原文} &
07-19 16:37 & 人工智能 & 1,670 \\
\rowcolor{paper}
中国新闻网 & AI研学营，“注水”很普遍
\newline\href{https://www.chinanews.com.cn/edu/2026/07-14/10658982.shtml}{查看原文} &
07-14 13:21 & AI & 4,112 \\
光明网 & 首届北京市中学生人形机器人足球赛总决赛举行
\newline\href{https://edu.gmw.cn/2026-07/23/content_38903218.htm}{查看原文} &
07-23 & 机器人 & 1,690 \\
\rowcolor{paper}
商汤科技 & 全球首个端侧驱动本体的具身世界模型！大晓机器人发布 Kairos 3.1
\newline\href{https://www.sensetime.com/cn/news/51170767}{查看原文} &
07-28 & 人工智能、AI、机器人 & 3,991 \\
科技日报 & 全国首个具身智能发展局揭牌
\newline\href{https://www.stdaily.com/web/gdxw/2026-07/29/content_555090.html}{查看原文} &
07-29 10:08 & 具身智能 & 607 \\
\rowcolor{paper}
新华网 & AI奢护新范式
\newline\href{https://www.news.cn/tech/20260724/4162865b956b492f9f7b3ac15d1e8626/c.html}{查看原文} &
07-24 & AI & 87 \\
\bottomrule
\end{longtable}

上述 6 条候选均保存了清洗正文和详情快照。详情快照合计 309,371 Bytes；本轮列表与详情
网络响应合计 993,007 Bytes。标题命中词与规则版本已写入候选提取元数据，后续可解释
“为什么这篇文章被接受”。

\section{质量、安全与独立审查}

\subsection{最终质量门}

\noindent\begin{tabularx}{\textwidth}{L{4.0cm}Y}
\toprule
\rowcolor{teal}
\color{white}\bfseries 检查 & \color{white}\bfseries 结果 \\
\midrule
后端自动化测试 & 112 项全部通过，总覆盖率 86\%，标题相关性模块覆盖率 100\% \\
\rowcolor{paper}
代码质量 & Ruff 格式/lint 通过；strict mypy 检查 46 个源文件通过 \\
真实服务集成 & PostgreSQL、MinIO、迁移、快照、幂等、租约、重试与无关详情阻断通过 \\
\rowcolor{paper}
前端合同 & OpenAPI 漂移检查、Prettier、ESLint、TypeScript、Vitest 3/3、生产构建通过 \\
运行环境 & Doctor、pgvector 0.8.1、Alembic \texttt{20260729\_0003}、MinIO Bucket 通过 \\
\rowcolor{paper}
交付检查 & Compose 渲染、Shell 语法、\texttt{git diff --check}、项目凭据正则扫描通过 \\
\bottomrule
\end{tabularx}

\subsection{独立审查发现并修复的问题}

\begin{itemize}[leftmargin=2em,itemsep=0.25em]
  \item 收紧 \textit{agent}、BCI、UAS 等英文歧义，补充 DeepSeek、ChatGPT、GPT、chatbot、
        intelligent system 与 Unicode 短横线形式；
  \item 防止“深度学习”在普通教育语境误收，防止“智能体检、智能体育、智能体系”等复合词误收；
  \item 数据库降级前恢复旧活动来源版本，避免旧 Worker 将新版本误判为无过滤版本；
  \item 修复重试时过滤计数重复累加，终态只保存一次扫描值；
  \item 详情页第一次请求也执行单来源限速，避免列表请求后立即连续访问。
\end{itemize}

\subsection{安全边界}

系统只允许来源登记表中的 HTTPS 域名和路径，逐次校验 DNS 与重定向目标，拒绝私网、环回、
链路本地、元数据和 Fake-IP/基准地址。响应设置连接/读取/总超时、最大字节、内容类型和
单来源速率。日志不记录完整正文、Cookie、授权头或签名 URL。抓取 Worker 已在验收后停止，
当前未完成任务和来源租约均为 0。

\section{当前边界与下一阶段建议}

\subsection{当前明确边界}

\begin{itemize}[leftmargin=2em,itemsep=0.25em]
  \item 当前只抓取 8 个已审核来源，不做全网搜索，也不接受任意 URL；
  \item 当前规则为确定性标题门禁，可能有意漏掉标题没有明确 AI 信号的新术语；
  \item 当前交付为服务器可部署形态，正式主机的域名/TLS、认证、备份和外部监控另行建设；
  \item 本阶段不做摘要、聚类、选题评分、文案生成、生图或自动发布；
  \item Docker 镜像在正常网络带宽下的完整构建和 gitleaks 工具扫描仍属于部署工具跟进项。
\end{itemize}

\subsection{下一阶段：LangGraph 工作流}

建议下一阶段使用 LangGraph 编排需要状态和分支的 AI 工作流，但不替代本阶段的调度器、
Worker、数据库租约和安全抓取。建议节点为：

\begin{center}
\statuspill{ocean}{读取候选正文}
$\rightarrow$ \statuspill{ocean}{摘要与结构化分类}
$\rightarrow$ \statuspill{ocean}{事件聚类/去重}
$\rightarrow$ \statuspill{ocean}{选题资格与评分}
$\rightarrow$ \statuspill{teal}{人工审核/生成}
\end{center}

LangGraph 节点使用 candidate ID 读取已保存正文和快照；原文链接用于引用与人工核验，
不在每个节点重新访问官网。这样模型服务不可用或工作流重试时，也不会重复抓取权威网站。

\begin{focusbox}
\textbf{\sffamily\color{navy}阶段结论}\quad
第一部分已经形成可信、可持续的事实输入边界：来源固定、主题可解释、无关内容可拒绝、
正文与原始网页可追溯、失败和重试可恢复。系统已具备进入第二部分“数据治理、摘要与选题”
建设的前提条件。
\end{focusbox}

\end{document}
